\documentclass[11pt]{article}
\usepackage[margin=1in]{geometry}
\usepackage{amsmath,amssymb,amsthm}
\usepackage[T1]{fontenc}
\usepackage{hyperref}

\newcommand{\Lip}{\operatorname{Lip}}
\newcommand{\Free}{\mathcal F}
\newcommand{\pten}{\widehat{\otimes}_{\pi}}

\title{Status Note: Vector-Valued Lipschitz Diameter Questions}
\author{Run fa\_banach\_001}
\date{2026-06-26}

\begin{document}
\maketitle

\paragraph{Status.}
\texttt{literature\_already\_answered (scoped)}.  This note records later
literature answers to selected questions from
Becerra Guerrero--Lopez-Perez--Rueda Zoca, arXiv:1512.03558,
\emph{Octahedrality in Lipschitz-free Banach spaces} \cite{BGLRZ1512}.

\paragraph{Source questions.}
In Section 3 of \cite{BGLRZ1512}, the authors ask whether, under the
hypotheses of their Theorem 2.4, $\Lip(M,X^*)$ satisfies the strong diameter
two property.  The same section asks for assumptions on a metric space $M$ and
a non-zero Banach space $X$ ensuring that the vector-valued free space
$\Free(M,X)=\Free(M)\pten X$ has slice-D2P, D2P, or SD2P.  In later notation
these are cited as Questions 3.1 and 3.3 of \cite{BGLRZ1512}.

\paragraph{Supporting answer.}
Langemets--Rueda Zoca, arXiv:1905.09061, \emph{Octahedral norms in duals and
biduals of Lipschitz-free spaces} \cite{LRZ1905}, explicitly cites
\cite{BGLRZ1512}.  Theorem 2.4 of \cite{LRZ1905} proves that if the pair
$(M,X)$ has the contraction-extension property and $M$ is unbounded or is not
uniformly discrete, then $\Lip(M,X)$ has the symmetric strong diameter two
property.  The paper states immediately after the theorem that this improves
\cite[Theorem 2.4]{BGLRZ1512} and gives a positive answer to
\cite[Question 3.1]{BGLRZ1512}.

Theorem 3.1 of \cite{LRZ1905} proves that if $\Lip(M)$ is octahedral, then
$\Lip(M,X^*)$ is octahedral for every Banach space $X$.  Corollary 3.2 then
states that, for a metric space $M$, the following are equivalent:
\[
  \Free(M,X)\text{ has the SD2P for every Banach space }X,
\]
\[
  \Free(M)\text{ has the SD2P},
  \qquad
  \Free(M)\text{ has no strongly exposed point}.
\]
The supporting paper states that this corollary answers
\cite[Question 3.3]{BGLRZ1512}.

\paragraph{Scope limitation.}
This packet is not a new proof.  It records an already-known literature
answer.  The clean criterion above is for the SD2P component of the
vector-valued free-space question.  The source question also mentions slice-D2P
and D2P; this packet does not claim a full classification of those weaker
variants beyond consequences of the SD2P result.

\begin{thebibliography}{9}

\bibitem{BGLRZ1512}
J. Becerra Guerrero, G. Lopez-Perez, and A. Rueda Zoca,
\emph{Octahedrality in Lipschitz-free Banach spaces},
arXiv:1512.03558.

\bibitem{LRZ1905}
J. Langemets and A. Rueda Zoca,
\emph{Octahedral norms in duals and biduals of Lipschitz-free spaces},
arXiv:1905.09061.

\end{thebibliography}

\end{document}
